\answer
\begin{enumerate}

%% 2章 問1（2.2節）
\item[2.2節 問1]
アクセプタをドーピングしたのでp形半導体である。
多数キャリアは正孔で$p \approx N_a = 10^{17}\,\mathrm{cm^{-3}}$。
少数キャリアは電子で，質量作用の法則\ref{eq2.1}より
\begin{equation*}
n = \frac{n_i^2}{p} = \frac{(1.5\times10^{10})^2}{10^{17}}
\approx 2.3\times10^{3}\,\mathrm{cm^{-3}}
\end{equation*}

%% 2章 問2（2.3節）
\item[2.3節 問2]
$v_d = \mu\mathcal{E}$の右辺$\mu\mathcal{E}$の単位は
\begin{equation*}
\frac{\mathrm{m^2}}{\mathrm{V\cdot s}}\times\frac{\mathrm{V}}{\mathrm{m}} = \mathrm{m/s}
\end{equation*}
となり，左辺（速度）と一致する。
式\ref{eq2.2}の右辺$en\mu\mathcal{E}$の単位は
\begin{equation*}
\mathrm{C}\times\mathrm{m^{-3}}\times\frac{\mathrm{m^2}}{\mathrm{V\cdot s}}\times\frac{\mathrm{V}}{\mathrm{m}}
= \frac{\mathrm{C}}{\mathrm{s\cdot m^2}} = \mathrm{A/m^2}
\end{equation*}
となり，左辺（電流密度）と一致する。

%% 2章 問3（2.4節）
\item[2.4節 問3]
空乏層の外には電場がないから，
空乏層の中のn形側の正の固定電荷とp形側の負の固定電荷は，全体として等しい量でつり合っている。
固定電荷の量は「イオンの濃度$\times$空乏層に入る幅」で決まる。
濃度の低い側は単位体積あたりのイオンが少ないので，
相手側と同じ量の電荷を出すには広い幅が要る。
したがって空乏層は濃度の低い側に深く広がる
（両側の幅を$x_p$，$x_n$とすると$N_a x_p = N_d x_n$，つまり幅は濃度に反比例する）。

%% 2章 問4（2.5節）
\item[2.5節 問4]
n形半導体で$\phi_m = 5.1\,\mathrm{eV} > \phi_s = 4.3$\,eVだから，
電子は半導体から金属へ移って界面に空乏層ができ，
整流性のショットキー接合になる（\ref{tab2.1}）。

\end{enumerate}
